\begin{circuitikz}[scale=0.8]
    \centering
    \draw (0,-0.3) to[V] (0,4.2)
    -- (1,4.2) to[opening normal closed switch, l={Schalter wird bei $t_\mathrm{0} = 5\tau$ geöffnet}] (3,4.2)
    -- (4,4.2);
    \draw (4,4.2)   to [C,i,v, name=C, l={$C$}] (4,2.2);
    \draw (4,0.2) -- (4,-0.2);
    \draw (4,-0.3) -- (0,-0.3);
    \draw (4,2.2)   to [R,v, name=R, l={$R$}] (4,-0.3);
    \draw[-latex, thick, color=cyan] (-0.7,2.55)  to node[midway,left, color=cyan] {$U_\mathrm{q}$}
    (-0.7,1.35);
    \iarrmore{C}{$i(t)$};
    \varrmore{C}{$u_\mathrm{C}(t)$};
    \varrmore{R}{$u_\mathrm{R}(t)$};

\end{circuitikz}

\onslide<2->{
    \begin{tikzpicture}
    \begin{axis} [
            xlabel={Zeit},
            ylabel={\textcolor{voltage} {$u_\mathrm{C}(t)$}\only<3->{/\textcolor{current}{$i(t)$}} } ,
            % ylabel={
            % 	\textcolor{cyan}{$U_\mathrm{q}$}}
            % 	\only<3->{/ \textcolor{voltage}{$u_\mathrm{C}(t)$}}
            % 	\only<4->{/ \textcolor{current}{$i(t)$}},
            xmin=0, xmax=12.5,
            ymin=-0.2, ymax=1.2,
            ytick={0, 0.37, 0.63, 1},
            yticklabels={%
                    0,%
                    37\%,%
                    63\%,%
                    100\%
                },
            %yticklabels={0,37\%, 63\%, 100\%},
            xtick={0, 1, 2, 3, 4, 5, 6, 7, 8, 9, 10, 11, 12},
            xticklabels={0, \(\tau\), \(2\tau\), \(3\tau\), \(4\tau\), \(5\tau\), \(6\tau\), \(7\tau\), \(8\tau\), \(9\tau\), \(10\tau\), \(11\tau\), \(12\tau\)},
            width=10cm,
            height=3.5cm,
            axis lines=middle,
            every axis x label/.style={at={(current axis.right of origin)},anchor=west},
            every axis y label/.style={at={(current axis.above origin)},anchor=south},
            font=\small,
            legend style={font=\scriptsize}
        ]
        \addplot [voltage, domain=0:12, samples=100] {ifthenelse(x<0, 0, 1 - exp(-(x)))};
        \only<3->{\addplot [current, domain=0:12, samples=100] {ifthenelse(x<0, 0,     exp(-(x)))};}

        %Schalter markierung in Graph
        \addplot [cyan,thick, dotted] coordinates {(5,1) (5,0)};

        % Markierung der Zeitkonstanten
        \draw[dashed] (axis cs:1,0) -- (axis cs:1,0.63) node[right] {};
        \draw[dashed] (axis cs:0,0.63) -- (axis cs:1,0.63) node[right] {};
        
        \only<3->{\draw[dashed] (axis cs:1,0) -- (axis cs:1,0.37) node[right] {};
            \draw[dashed] (axis cs:0,0.37) -- (axis cs:1,0.37) node[right] {};}
        \node[right] at (5,0.5) {$t_\mathrm{0}=5\tau$};

        % Punkte zur Hervorhebung
        \addplot[
            only marks,
            mark=*,
            mark options={scale=0.7, fill=black}
        ]  coordinates {
                (1, 0.63)

            };
        \only<3->{ \addplot[
                only marks,
                mark=*,
                mark options={scale=0.7, fill=black}
            ]  coordinates {
                    (1, 0.37)

                };}

    \end{axis}
\end{tikzpicture}
}